\title{\textbf{Slow Wave}: A Mechanism-Decomposed Test Bench for
Sleep-Inspired \emph{Memory Consolidation} in \emph{Continual-Learning}
\emph{LLM Agents}}

\author{The Slow Wave Project}
\date{\today}
\maketitle

\begin{abstract}
Long-running, frozen-weights large language model (LLM) agents ingest far more
data than ever affects mission performance. Biology faces an analogous problem
and answers it with \emph{sleep}: an offline cycle that replays recent
experience, consolidates episodic traces into durable semantic memory, and
downscales low-value synapses while protecting salient ones. We ask whether an
engineered analogue of this \emph{sleep/dream} cycle improves long-horizon
\emph{continual learning} for an \emph{LLM agent} whose weights are frozen ---
i.e.\ whether \emph{memory consolidation}, not weight update, can carry the
effect. We build \textbf{Slow Wave}, an instrumented, reproducible test bench:
a synthetic continual task stream with injected distractors and
\emph{ground-truth relevance labels}, a dual-store (episodic/semantic/archival)
memory substrate, and a four-operator dream engine (replay, transfer,
downscale, generative-augment) whose operators are independently ablatable. A
preregistered, budget-controlled control battery of \SWnArms{} arms --- including a
ground-truth-blind random-pruning negative control, an oracle prune ceiling, a
long-context baseline, and an A/A noise floor --- is run over an
arm\,$\times$\,regime\,$\times$\,seed grid (\SWnCells{} cells,
\SWnSeeds{} seeds). Beyond downstream accuracy we measure consolidation
\emph{quality} directly: precision/recall/F1 of discarding distractors versus
retaining signal. \textbf{All reported numbers are a mechanism demonstration
produced under a deterministic mock LLM in the synthetic stream, and are not
claims about a real Claude model.} In this regime the full dream cycle improves
final average accuracy over the no-sleep baseline on the preregistered primary
endpoint (paired $\Delta=\SWprimaryDiff$, 95\% CI
$[\SWprimaryCIlo,\SWprimaryCIhi]$), with a registered secondary contrast confirming
it also beats a replay-only control; because the realised token and memory
budgets are not matched within tolerance, we report these as cost-effectiveness
results on the accuracy-vs-compute Pareto frontier rather than matched-budget
wins. A replay-targeting analogue moves signal into durable memory as designed,
though its standardised margin over a sleep-targeted-memory-reactivation
benchmark is inflated by the mock's determinism and is not a meaningful
comparison. Reported with equal prominence as first-class findings, both an
uncapped long-context baseline and a simpler deployable shallow-synthesis control
(\texttt{reflection}) dominate the full cycle on the accuracy-vs-token frontier ---
long-context on raw accuracy at maximum memory cost (with no cost-adjusted crossover in
the swept stream lengths), reflection at higher accuracy and lower token cost. We
treat negative and cost-unfavourable patterns with the same rigor as the
positive result, and release the apparatus, preregistration, manifests, and a
one-command reproduction of every figure and number.
\end{abstract}

\noindent\textbf{Keywords:} memory consolidation; continual learning; LLM
agents; sleep/dream; experience replay; catastrophic forgetting; synaptic
homeostasis; preregistration; reproducibility.
